\documentclass{aa}  
%\documentclass[referee]{aa}    % for a referee version

\usepackage{graphicx}
\usepackage{txfonts}
\usepackage{subcaption}         % necessary for continued figures, example in section 3
                                % and appendix
\usepackage{lscape}             % to rotate a single page table, example in appendix.
                                % For landscape tables, see the longtable examples.
\usepackage{placeins}           % useful with \FloatBarrier, to keep 
                                % onecolumn floats from drifting to the next section
                                

%%%%%%%%%%%%%%%%%%%%%%%%%%%%%%%%%%%%%%%%
\usepackage[]{hyperref}
% To add links in your PDF file, use the package "hyperref"
% with options according to your LaTeX or PDFLaTeX drivers.
%%%%%%%%%%%%%%%%%%%%%%%%%%%%%%%%%%%%%%%%

\newcommand{\Gaia}{\textit{Gaia}}
\newcommand{\bootes}{Bo\"otes III}
\newcommand{\boo}{Boo III}
% mathematics commands
\newcommand{\rmax}{\ensuremath{r_\textrm{math}}}
\newcommand{\vmax}{\ensuremath{\mathrm{v}_\textrm{math}}}
\newcommand{\pmra}{\ensuremath{\mu_{\alpha*}}}
\newcommand{\pmdec}{\ensuremath{\mu_\delta}}
\newcommand{\masyr}{\ensuremath{\textrm{mas yr}^{-1}}}
\newcommand{\kms}{\ensuremath{\textrm{km s}^{-1}}}
\newcommand{\vlos}{\ensuremath{\mathrm{v}_\text{los}}}
\newcommand{\kpc}{\ensuremath{\textrm{kpc}}}
\newcommand{\Gyr}{\ensuremath{\textrm{Gyr}}}
\newcommand{\Mo}{\ensuremath{\mathrm{M}_\sun}}


\begin{document}

%%%%%%%%%%%%%%%%%%%%%%%%%%%%%%%%%%%%%%%%
% if you use custom commands in your title,
% ensure to check your title when submitting!
%%%%%%%%%%%%%%%%%%%%%%%%%%%%%%%%%%%%%%%%


   \title{Where is Bo\"otes III's tidal stream?}

   %\subtitle{}

%%%%%%%%%%%%%%%%%%%%%%%%%%%%%%%%%%%%%%%%
% Please separate each author with the \and command
%
% Use the \corrauth to provide the corresponding
% author address. It will be automatically inserted as 
% footnote in the PDF output.
%
% Please DO NOT include ORCIDs next to author names.
% Instead, please provide an active address for each coauthor:
% it will be automatically extracted by EDPS editorial system, 
% and co-authors will be be able to authenticate their ORCID.
%
% Only authenticated ORCIDs will be taken into account.
% ORCIDs included here will be removed.
%%%%%%%%%%%%%%%%%%%%%%%%%%%%%%%%%%%%%%%%

   \author{D. A. Boyea\inst{1}\corrauth{danielaboyea@gmail.com}        % use \corrauth for the corresponding author
       \and all
        %\and F. Déliat\inst{1}\email{fab@nestor-edp.org}\fnmsep\thanks{NASA fellow (shows the usage of elements in the author field)}
        }

   \institute{UVic, V8W 2Y2, BC, Canada
       %\and
   }

   \date{Received date}

% \abstract{}{}{}{}{}
% 5 {} token are mandatory
 
  \abstract
  % context heading (optional)
  % {} leave it empty if necessary  
   {}
  % aims heading (mandatory)
   {We aim to update the observational properties of \bootes{} with modern datasets and simulate its tidal evolution to understand the formation of a possible stream.} 
  % methods heading (mandatory)
   {We select stars from \textit{Gaia} to derive structural properties using standard Monte Carlo Markov Chain fitting. We then use N-body collisionless dark matter simulations to predict the tidal evolution of \bootes{} in the Milky Way potential.
   }
  % results heading (mandatory)
   {We present some of the first structural parameter derivations in over a decade, improving on past determinations. Our orbtial analysis indicates that \boo{}'s orbit depends heavily on the current heliocentric distance, allowing pericentrew from between 1.5 and 18 kpc (2-$\sigma$ confidence range). With full N-body simulations, we dermine that \bootes{} requires a larger pericentre of at least 18 kpc in order to reproduce the presenty observed properties. Additionally, while some models lose more than half their initial stellar material, we find that all models form tidal streams which are diffuse (spanning $\sim 5$ degrees in width) and orders of magnitude fainter than the currently detectible in the \textit{Gaia} or DELVE (Dark Energe camera Local Volumne Experiement) surveys. Excellent photometric depth and discrimination of background sources will be required to observe such features. 
   }
  % conclusions heading (optional), leave it empty if necessary
   {
       If the Styx stream exists, its current apparent properties are inconsistent with an origin from \boo{}.  More precise determinations of \boo{}'s distance will be essential to understanding its dynamical evolution and consistency with $\Lambda$CDM cosmology. 
   }

   \keywords{
       dSph galaxies --
       tidal effects
               }

   \maketitle


%%%%%%%%%%%%%%%%%%%%%%%%%%%%%%%%%%%%%%%%%%%%%%%%%%%%%%%%%%%%%%
\section{Introduction}

Of the ${\sim}60$ dwarf galaxies known to orbit the Milky Way, only five show obvious signs of tidal perturbations. The most dramatic examples are the Sagittarius dwarf spheroidal (dSph), forming streams wrapping across the entire sky, and the Large Magellanic Cloud's extensive gas tail (LMC). Otherwise, only Crater II, Antlia II, and Tucana III have detected streams. Indeed, most dwarf galaxies have central densities too high for tidal affects to perturb their stellar component \citep[e.g.,][]{pace+2020}.

\bootes{} (\boo{}) has always thought to be in the throws of tidal disruption, forming the extensive Styx stellar stream. \boo{} was discovered with Styx in \citet{grillmair2009} through a matched filter method (where stars are weighted based on their CMD position to increase the contrast of a specific stellar population). Along with an initially high velocity dispersion of $14$ km\,s, later proper motions suggesting a small pericentre, and the apparent consistency of \boo{}'s orbit with Styx lead to a clear conclusion: \boo{} must be undergoing active tidal disruption and is the proginiter of Styx. 

However, \citep{jensen2026} find that Styx does not appear to be reflected in any overdensities in Blue Horizontal Branch or Red Giant Branch stars, nearly no \Gaia{} stars are consistent with being Styx members, and only a tenous detection in matched filter analysis can be confirmed. A proper dynamical model of \boo{} is needed to understand if a stream like Styx plausably formed from \boo{} and whether \boo{} may reveal something about tidal evolution in the Milky Way. 



\section{The observational status of Bootes III}
In this section, we review the current observational properties of \boo{} and derive new structural properties, kinematic information, and metallicity dispersions. We also outline the datasets used later for comparison with our models.

Table~\ref{tab:obs_props} summarizes our recommended observational properties of \boo{}.

\begin{table}
    \begin{tabular}{l l l l}
        \hline\hline
        parameter & units & value & reference \\
        \hline
        $\alpha$ & deg & $209.47 \pm 0.11$ & section~\ref{sec:struct_params} \\
        $\delta$ & deg & $26.68 \pm 0.06$ & section~\ref{sec:struct_params}\\
        DM & mag & $18.34 \pm 0.19$ & \citet{vivas+2020} \\
        distance & kpc & $46.56 \pm 4$ & \citet{vivas+2020} \\
        $\pmra$ & $\masyr$ & $-1.16 \pm 0.02 \pm 0.017$ & \citet{battaglia+2022} \\
        $\pmdec$ & $\masyr$ & $-0.88  \pm 0.01 \pm 0.017$ & \citet{battaglia+2022} \\
        $\vlos$ & $\kms$ & ** & (Ting Li in prep) \\
        $\sigma_\mathrm{v}$ & $\kms$ & ** & (Ting Li in prep) \\
        $R_h$ & arcmin & $44_{-6}^{+7}$ & section~\ref{sec:struct_params}\\
        ellipticity & --- & $0.42_{-14}^{+0.11}$ & section~\ref{sec:struct_params}\\
        position angle & deg & $89 \pm 9$ & section~\ref{sec:struct_params}\\
        $M_V$ & mag & $-5.1 \pm 0.24$ & section~\ref{sec:luminosity} \\
        $M_\star$ & $\Mo$ & $(1.4 \pm 0.3) \times 10^4$ & section~\ref{sec:luminosity} \\
        %[Fe/H] & dex & $-2.5 \pm 0.1$ & section \\
        age & Gyr & 12: & assumed \\
        \hline
    \end{tabular}

    \caption{The observational properties of \boo{} from the literature and derived here.
        Rows: right ascension ($\alpha$), declination $\delta$, distance modulus, distance, absolute proper motion in right ascension $\mu_\alpha \cos\delta$, proper motion in declination $\mu_\delta$, line of sight (los) velocity $\mathrm{v_{los}}$, LOS velocity dispersion $\sigma_\textrm{v}$, tentative half-light radius $R_h$, ellipticity, position angle, absolute V-band magnitude, total stellar mass, and metallicity.
    }
    \label{tab:obs_props}
\end{table}

\subsection{\textit{Gaia} observations\label{sec:gaia_obs}}



\begin{figure*}
    \centering
    \includegraphics{figures/gaia_tangent_cmd_pm.pdf}
    \caption{The distribution of \Gaia{} stars in the tangent plane (left),  the \Gaia{} colour-magnitude diagram (right), and proper motion space (lower right). The small grey stars satisfy basic quality cuts, the larger solid blue points are likely members based on CMD and PM likelihoods (see text), and the red diamonds are additionally RV members.}

      \label{fig:gaia_tangent}
\end{figure*}




\subsection{Structural Parameters\label{sec:struct_params}}

\begin{figure}
    \centering
    \includegraphics{figures/mcmc_density_profiles.pdf}
    \caption{
        The observed density profile of \boo{} plotted as log surface density $\Sigma_\star$ agains log elliptical radius $R_\textrm{ell}$ . 
        The blue points with error bars are the binned density profile, whereas the orange and green transparent lines are individual samples from the Plummer and S\'ersic MCMC structural fits to \Gaia{} data, respectively.
        \newline{}
        \boo{} is characterized by a diffuse and extended surface density profile. Because of \boo{}'s low surface brightness, the background/forground limit our ability to determining the density profile.
    }
    \label{fig:density_profile_fit}
\end{figure}




\begin{figure}
    \centering
    \includegraphics{figures/RV_hist_mcmc.pdf}
    \caption{
        Our MCMC fits to the velocity profile of \boo{}. 
        Black lines show the individual MCMC Gaussian samples, and the orange points with error bars show the 16 los-velocity member stars with their associated uncertainties (placed arbitrarily).
    }
    \label{fig:sigma_v_fit}
\end{figure}



\section{Numerical methods}

\subsection{Selection of initial dark matter halos}

\begin{figure*}
    \centering
    \includegraphics{figures/initial_halo_selection.pdf}
    \caption{
        The initial $\rmax$ and $\vmax$ which \citep{EN2021} predicts will acchieve \boo{}'s velocity dispersion (curves), and the cosmologically expectations for \boo{}'s vmax given its stellar mass (green horizontal line), and \boo{}'s rmax given a vmax (grey slanted line). 
        \textbf{Left:} the required halos matching the final velocity dispersion for different orbits. 
        \textbf{Middle:} The required halos for different final velocity dispersions on the fiducial 12kpc orbit. 
        \textbf{Right:} The required halos for different lengths of orbit on the fiducial 12kpc orbit. 
        \newline
        \textbf{Explanation:} \boo{} requires either an unlikely orbit (e.g. 26kpc pericentre or a low number of pericentres) or an unexpectedly massive and concentrated initial halo.
    }
    \label{fig:initial_halo_selection}
\end{figure*}



\begin{figure}
    \centering
    \includegraphics{figures/pericentre_hist_corr.pdf}
    \caption{
        The distribution of possible pericentres. 
    }
    \label{fig:peri_dist}
\end{figure}


\begin{figure*}
    \centering
    \includegraphics{figures/orbits.pdf}

    \caption{
        Six example orbits of \boo{} shown on the $x$--$y$ Galactocentric plane (left), the $y$--$z$ plane (middle), and in Galactocentric radius versus time (right).
        \newline{}
        \textbf{Explanation:} A wide range of orbits is possible for \boo{}. Orbits with smaller pericentres also have smaller apocentres and shorter orbital periods. 
    }
    \label{fig:orbits}
\end{figure*}


%%%%%%%%%%%%%%%%%%%%%%%%%%%%%%%%%%%%%%%%%%%%%%%%%%%%%%%%%%%%%%%%%%%%%%%%%%%%%%%%
\section{Results}

\subsection{Tidal evolution}

\begin{figure}
    \centering
    \includegraphics[width=\columnwidth]{figures/fiducial_zoom_image.png}
    \caption{
        The final distribution of dark matter (purple) and stars (white) for the fiducial model in the $y$--$z$ Galactocentric plane. The density scale spans 5 orders of magnitude and is logarithmic.
        \newline
        \textbf{Explanation:} The model forms substantial streams. Stars span nearly a half-orbital period forward and backwards in the orbit. \boo{} has a radial enough orbit such that the stream forms shells at apocentres.
    }
    \label{fig:fiducial_sim_image}
\end{figure}


\begin{figure}
    \centering
    \includegraphics{figures/boundmass_sigma_v.pdf}

    \caption{
        The evolution of stellar velocity dispersion (left,  $\sigma_\text{los}$) and stellar/dark matter bound mass (right) over time for each simulation.
        \newline{}
        \textbf{Explanation:} The observed stellar velocity dispersion is acchieved, but the total stellar mass loss varies depending on the model. \boo{} may have lost a majority of its initial stellar mass!
    }

\end{figure}


\begin{figure*}
    \centering
    \includegraphics{figures/final_stellar_streams.pdf}
    \caption{
        The distribution of stars onsky at the end of each simulation. 
        \newline
        \textbf{Explanation:}
        Stream brightness and morphology depends most strongly on the pericentre and number of pericentres. Smaller pericentres form more significant and thicker streams (compare fiducial, 7kpc, and 26kpc pericentres). However, experiencing only two small pericentres only formss a meager tidal stream. 
        Across models, the tidal stream is relatively faint (surface densities of $\lesssim 1 \text{stars\ deg}^{-2}$. 
    }
    \label{fig:predicted_stars_onsky}
\end{figure*}




\begin{figure}
    \centering
    \includegraphics{figures/along_stream_density.pdf}
    \caption{
        The predicted surface brightness along the stream axis (marked in Fig.~\ref{fig:predicted_stars_onsky}).
        \newline{}
        \textbf{Explanation:} Within the inner 4 degrees, the surface density drops but with a possible halo-like extension. Outside of the inner galaxy, the stream remains extemely diffuse, with less than 0.1 stars per square degree predicted in \textit{Gaia} data. \boo{}'s stream remains unacessable without extremely deep, low-foreground data. 
    }
    \label{fig:stream_brightness}
\end{figure}

\subsection{Prediction of a stream}


%%%%%%%%%%%%%%%%%%%%%%%%%%%%%%%%%%%%%%%%%%%%%%%%%%%%%%%%%%%%%%%%%%%%%%%%%%%%%%%%
\section{Conclusions}


%%%%%%%%%%%%%%%%%%%%%%%%%%%%%%%%%%%%%%%%%%%%%%%%%%%%%%%%%%%%%%%%%%%%%%%%%%%%%%%%
\begin{acknowledgements}
\end{acknowledgements}



\bibliographystyle{aa}
\bibliography{paper.bib}



\begin{appendix}


\section{Extended structure search in \textit{Gaia}}

To search for structure over an extended area in \textit{Gaia}, we download and reprocess \textit{Gaia} sources following \citet{MV2020}. For simplicity, we select all \textit{Gaia} sources within a rectangle of 10 degrees by 10 degrees in RA and Dec of the literature centre of \boo{} (209.3, 26.8). Next, we apply the quality cuts from \citet{MV2020, MV2020a}
\begin{itemize}
    %\item All relevant parameters are defined ($\alpha$, $\delta$, $\pmra$, $\pmdec$, $\sigma_\pmra$, $\sigma_\pmdec$, $G$, $BP$, $RP$, $\sigma_G$, $\sigma_{BP}$, $\sigma_{RP}$, $\mu$, $\delta \mu$).
    %\item Proper motions are between $-10$ and $10$ $\masyr$
    \item The parallax is within $3\sigma$ of \boo{}'s parallax (includeing \boo{}'s parallax uncertainty, the sources parallax uncertainty, and the systemic parallax uncertainty from CITE. 
    \item The colour flux excess is within $3\sigma$ of the expected value from XXX.
    \item BP - RP colour is between XX and XX
    \item 
\end{itemize}

\section{Structural properties in DELVE}

For DELVE (DR2), we select all stars within 6 arcminutes of BooIII's literature centre (209.3, 26.8).\footnote{DELVE DR3 has incomplete coverage near Boo3.} We require the sources to further satisfy:  
\begin{itemize}
    \item g and r PSF magnitudes and their errors are defined
    \item g and r PSF magnitudes greater than the 5-$\sigma$ PSF depth (24.3, and 23.9, respectively, \citep{delve_dr2})
  - likely star source ($\texttt{extended\_class\_g} \leq 1$)
\end{itemize}

In total, this returns 558,046 sources. We then correct the magnitudes for dust extinction using the @SFD dustmap with coefficients from @aboot+2018 with the dustmaps package.

We believe that DELVE has slightly inhomogenous completeness near \boo{}. Fig.~\ref{fig:delve_completeness} illustrates our derived 5-sigma point magnitude depth and the matched filter density as a function of position. We calculate the 5-sigma point source depth by taking the median magnitude of stars with magnitude uncertainties in the g band between (0.2 and ???), when spatially binning in the tangent plane. Patterns in the 5-sigma point source depth are refelcted in the matched filter density. We conclude that a proper correction of the inhomogenous completeness of DELVE is necessary to properly search for faint, extended structure. 



\begin{figure}
    %\includegraphics{figures/temp}
    \caption{Our derived matched filter density for DELVE data. 
    }
\end{figure}

For simplicity, we instead reduce our magnitude limit to $g=22$. While this appears to be extreme, we avoid possible biases from improper 

\section{Cored dark matter haloes}


\begin{figure*}
   \includegraphics{figures/cored_tracks.pdf}
\end{figure*}

\end{appendix}
\end{document}
